% =====================================================================
%  Section 4: Experiments
%  Main setting: LLaMA-2-7B full fine-tuning, Alpaca-GPT4,
%  10% selection, nine evaluation tasks across five capability groups.
% =====================================================================
\section{Experiments}
\label{sec:experiments}

\begin{table*}[t]
\centering
\begin{threeparttable}
\caption{\textbf{Performance comparison of baselines using 10\% of
the IT data.} This table shows the results of various tasks across
multiple baselines. \textbf{Bold} and \underline{underline} represent the best and 
second-best performance respectively in each column 
\emph{among 10\% selection methods (R10, ID~02--09);} ID~00 and 
ID~01 are shown as lower/upper-bound references and excluded 
from the ranking. $\Delta\,(\uparrow)$ indicates the performance improvement
relative to ID~01.}
\label{tab:main-results}
\small
\setlength{\tabcolsep}{2.9pt}
\renewcommand{\arraystretch}{1.01}

\begin{tabular}{@{}llcccccccccc@{}}
\toprule
 & &
 \multicolumn{2}{c}{Knowledge \& reasoning} &
 \multicolumn{2}{c}{Math} &
 \multicolumn{2}{c}{Code} &
 \multicolumn{2}{c}{Multilingual QA} &
 \multicolumn{2}{c}{Overall} \\
\cmidrule(lr){3-4}
\cmidrule(lr){5-6}
\cmidrule(lr){7-8}
\cmidrule(lr){9-10}
\cmidrule(lr){11-12}
ID & Method
& \shortstack[c]{MMLU\\[-1pt]{\scriptsize exam}}
& \shortstack[c]{BBH\\[-1pt]{\scriptsize hard tasks}}
& \shortstack[c]{SVAMP\\[-1pt]{\scriptsize robust math}}
& \shortstack[c]{GSM8K\\[-1pt]{\scriptsize grade math}}
& \shortstack[c]{MBPP\\[-1pt]{\scriptsize program}}
& \shortstack[c]{H-Eval\\[-1pt]{\scriptsize function}}
& \shortstack[c]{TydiQA\\[-1pt]{\scriptsize 9-lang F1}}
& \shortstack[c]{XQuAD\\[-1pt]{\scriptsize 12-lang F1}}
& \shortstack[c]{AVG\\[-1pt]{\scriptsize 8-task mean}}
& \shortstack[c]{$\Delta$\\[-1pt]{\scriptsize vs Full FT}} \\
\midrule

\rowcolor{gray!15}
\multicolumn{12}{@{}l}{\textbf{Data-side references}}\\
00 & Pure LLaMA-2-7B (No FT)
    & 45.82 & 36.35
    & 36.67 & 13.87
    & 0.00$^\dagger$ & 2.38$^\dagger$
    & 61.70 & 55.08
    & 31.48 & $-20.96\%$ \\
01 & Alpaca-GPT4 (Full FT)
    & 47.86 & 37.19
    & 40.33 & 16.60
    & 37.35 & 25.60
    & 60.74 & 52.99
    & 39.83 & --- \\
R10 & Random ($10\%$)
    & 45.04 & 35.54
    & 36.80 & 13.57
    & 35.41 & 24.51
    & \underline{60.83} & 52.28
    & 38.00 & $-4.61\%$ \\
\midrule
\rowcolor{gray!15}
\multicolumn{12}{@{}l}{\textbf{Static 10\% IT-data selection}}\\
02 & LIMA
    & 46.44 & 36.68
    & 36.20 & 16.38
    & 33.07 & 22.09
    & 59.28 & 51.14
    & 37.66 & $-5.45\%$ \\
03 & AlpaGasus
    & 46.17 & 36.36
    & 38.33 & 16.22
    & 37.50 & 23.94
    & 60.81 & \underline{53.87}
    & 39.15 & $-1.71\%$ \\
04 & Q2Q
    & 46.18 & 36.84
    & 34.33 & \underline{17.82}
    & 36.19 & 25.19
    & 60.42 & 52.34
    & 38.66 & $-2.93\%$ \\
05 & SelectIT
    & 46.02 & 38.37
    & 36.33 & 16.98
    & 36.58 & \underline{28.17}
    & 60.44 & \textbf{54.12}
    & 39.63 & $-0.52\%$ \\
06 & NAIT (All)
    & 46.08 & 40.02
    & \underline{39.67} & 17.06
    & 35.41 & 26.44
    & 60.03 & 52.87
    & 39.70 & $-0.34\%$ \\
\midrule
\rowcolor{gray!15}
\multicolumn{12}{@{}l}{\textbf{Dynamic 10\% selection}}\\
07 & \textsc{Data Agent (PPO)}
    & \underline{46.49} & 38.37
    & 37.67 & 16.48
    & \textbf{39.69} & \textbf{28.61}
    & 60.44 & 50.85
    & \underline{39.83} & \underline{$-0.02\%$} \\
08 & CO ($\lambda=0$)
    & 45.90 & \underline{40.21}
    & 36.67 & 16.53
    & 34.63 & 24.38
    & 60.12 & 50.70
    & 38.64 & $-2.99\%$ \\
\rowcolor{red!8}
09 & \textbf{\textsc{TAG}} ($\lambda=1$)
    & \textbf{47.71} & \textbf{40.36}
    & \textbf{40.33} & \textbf{18.27}
    & \underline{38.52} & 25.87
    & \textbf{61.13} & 53.43
    & \textbf{40.70} & \textbf{$+2.18\%$} \\
\bottomrule
\end{tabular}
\begin{tablenotes}[flushleft]
\footnotesize
\item[$\dagger$] MBPP grades pass@1 by executing the generated
function; the non-instruction-tuned base model emits free-form
text instead of parseable code, yielding $0$. HumanEval is a
signature+docstring completion task closer to pretraining, hence
the near-zero $2.38$.
\end{tablenotes}
\end{threeparttable}
\end{table*}

We evaluate \textsc{TAG} as a dynamic re-ranking method for
instruction-tuning data selection. Our primary setting is
full-parameter fine-tuning of LLaMA-2-7B~\cite{touvron2023llama2} on
Alpaca-GPT4~\cite{wang2023selfinstruct} with a fixed selection budget
of \(\rho{=}10\%\). Robustness experiments additionally vary the
backbone, using Qwen2.5-0.5B, Qwen2.5-7B~\cite{qwen2025technical},
and Mistral-7B~\cite{jiang2023mistral}, and the source dataset, using
WizardLM Evol-Instruct~\cite{xu2024wizardlm}. Within each
(backbone, dataset) setting, all selectors share the same training
hyperparameters and evaluation protocol; only the selected subset
differs. Full FT uses the full pool, while all selection methods use
\(10\%\), a budget chosen because random ratio sweeps show that
increasing random subsets up to \(50\%\) provides little additional
headroom.

We evaluate eight benchmarks spanning four capability families:
knowledge and hard reasoning (MMLU, BBH), mathematical reasoning
(SVAMP, GSM8K), code generation (MBPP, HumanEval), and multilingual QA
(TyDiQA, XQuAD). We report per-task scores, AVG, and \(\Delta\) over
Full FT; benchmark scoring conventions are summarized in the appendix.

\paragraph{Baselines.}
We compare against Full FT, random subsets, static instruction-data
selectors (LIMA, AlpaGasus, Q2Q, SelectIT, NAIT), the PPO-based Data
Agent, and CO, the Carrier-Only variant obtained by setting
\(\lambda{=}0\). CO is the direct baseline for isolating the
contribution of representation-anchor modulation.

\paragraph{Research questions.}
The experiments ask whether \textsc{TAG} (i) outperforms static and
dynamic selectors at the same \(10\%\) budget, (ii) remains robust
across backbones and source datasets, (iii) has a stable
anchor-strength choice, and (iv) adds only modest re-ranking overhead.


% \subsection{Main result}
% \label{sec:main-result}

% Table~\ref{tab:main-results} reports the main comparison under the same
% selection budget and evaluation protocol. The controlled comparison
% between the carrier-only baseline (CO) and \textsc{TAG} isolates the
% effect of representation-anchor modulation: both use the same dynamic re-ranking
% pipeline and carrier reward, while \textsc{TAG} additionally enables
% the anchor term. The key quantitative result is that \textsc{TAG}
% reaches an 8-task AVG of 40.70, the highest overall score among all
% 10\% selection methods.

% \paragraph{Effect of representation-anchor modulation.}
% \textsc{TAG} consistently improves over the CO baseline across the
% benchmark suite. Quantitatively, the anchor raises AVG from 38.64 to
% 40.70 (+2.06 pp), with especially clear gains on MBPP (+3.89 pp),
% SVAMP (+3.66 pp), and XQuAD (+2.73 pp). The gains span reasoning,
% math, code, and multilingual QA, suggesting that the anchor does not
% merely specialize selection to one capability family. This supports
% the central claim that a state-conditioned structural signal can
% improve dynamic data selection without seed examples, validation
% labels, or auxiliary scorers.

% \paragraph{Comparison to baselines.}
% Among methods using the same selection budget, \textsc{TAG} achieves
% the strongest overall performance, outperforming static selectors and
% the PPO-based Data Agent. On the same 10\% budget, this corresponds
% to a +0.87 pp gain over the PPO-based Data Agent and a +1.00 pp gain
% over the strongest static selector, NAIT. Static selectors rely on fixed
% or precomputed signals, whereas \textsc{TAG} refreshes a
% model-intrinsic representation signal during training. The random subset
% results serve only as budget controls, separating selection quality
% from the effect of using more data.

% \paragraph{Selection can beat scale.}
% The results show that a carefully selected subset can outperform full
% fine-tuning. In particular, \textsc{TAG} reaches 40.70 with only 10\%
% of the instruction-tuning data, compared with 39.83 for full
% fine-tuning, yielding a +0.87 pp absolute gain and a +2.18\% relative
% improvement. This highlights the value of adaptive selection criteria
% that are not only accurate, but also stable, reference-free, and
% computationally practical.


\subsection{Main result}
\label{sec:main-result}

Table~\ref{tab:main-results} reports the main comparison under the same
selection budget and evaluation protocol. The controlled comparison
between the carrier-only baseline (CO) and \textsc{TAG} isolates
representation-anchor modulation: both use the same dynamic re-ranking
pipeline and carrier reward, while \textsc{TAG} additionally enables
the anchor-projection factor. \textsc{TAG} reaches \textbf{40.70 AVG}, the highest
score among 10\% selection methods.

\paragraph{Effect of representation-anchor modulation.}
\textsc{TAG} consistently improves over the CO baseline across
benchmarks. The anchor raises AVG from 38.64 to 40.70
(+2.06 pp), with clear gains on MBPP (+3.89 pp), SVAMP (+3.66 pp),
and XQuAD (+2.73 pp). The gains span reasoning, math, code, and
multilingual QA, suggesting that the anchor does not merely
specialize to one capability family. This supports our claim that a
state-conditioned representation signal can improve dynamic data
selection without seed examples or auxiliary scorers. Among methods
using the same budget, \textsc{TAG}
achieves the strongest overall performance, outperforming static
selectors and the PPO-based Data Agent. At the same 10\% budget,
\textsc{TAG} gains +0.87 pp over Data Agent and +1.00 pp over
NAIT. Static selectors rely on fixed or precomputed signals, whereas
\textsc{TAG} recomputes its representation anchor from the current
model state at each refresh. The random subset results serve only as
budget controls,
separating selection quality from data scale.

\paragraph{Selection can beat scale.}
The results show that a carefully selected subset can outperform full
fine-tuning. With only 10\% data, \textsc{TAG} scores 40.70 versus
39.83 for Full FT, yielding +0.87 pp absolute and +2.18\% relative
gains. This highlights the value of adaptive selection criteria that
are not only accurate, but also stable, reference-free, and
computationally practical.


\begin{figure}[!ht]
\centering

% ============ TOP: compact Delta plot ============
\resizebox{0.97\columnwidth}{!}{%
\begin{tikzpicture}[every node/.style={font=\scriptsize}, >=Latex]

\path[use as bounding box] (-0.55,-0.95) rectangle (8.80,5.40);

% Y-axis grid  (왼쪽 숫자 0~4)
\foreach \y in {0,1,2,3,4} {
  \draw[gray!25, very thin] (0,\y) -- (7.2,\y);
  \node[anchor=east, font=\scriptsize] at (-0.05,\y) {\y};
}

% X-axis ticks (0, 0.1, 0.5, 1.0, 5.0)
\foreach \xpos/\lab in
  {0/0, 1.8/0.1, 3.6/0.5, 5.4/1.0, 7.2/5.0} {
  \draw[gray!50, very thin] (\xpos,0) -- (\xpos,-0.10);
  \node[anchor=north, font=\scriptsize] at (\xpos,-0.12) {$\lab$};
}

% Default lambda = 1.0 guideline
\draw[gray!60, dashed, thin] (5.4,0) -- (5.4,4.30);
\node[red!85!black, font=\small, anchor=south, align=center]
  at (5.4,4.35) {\textbf{default}\\[-1pt]$\lambda{=}1.0$};

% Per-task lines
\draw[blue!55, very thin] plot coordinates
  {(0,0) (1.8,0.26) (3.6,0.91) (5.4,1.81) (7.2,0.58)};
\draw[green!45!black!60, very thin] plot coordinates
  {(0,0) (1.8,0.66) (3.6,3.33) (5.4,3.66) (7.2,1.00)};
\draw[orange!75, very thin] plot coordinates
  {(0,0) (1.8,0.39) (3.6,2.72) (5.4,3.89) (7.2,1.56)};
\draw[violet!65, very thin] plot coordinates
  {(0,0) (1.8,0.44) (3.6,0.73) (5.4,1.01) (7.2,0.50)};

% Endpoint labels (좁은 구간에 모여있어 강제 분산)
\node[orange!85!black,font=\small, anchor=west] at (7.23,1.56) {MBPP};
\node[green!45!black, font=\small, anchor=west] at (7.23,1.05) {SVAMP};
\node[blue!55!black,  font=\small, anchor=west] at (7.23,0.55) {MMLU};
\node[violet,         font=\small, anchor=west] at (7.23,0.15) {TydiQA};

% Average line (굵게)
\draw[red!75!black, ultra thick] plot coordinates
  {(0,0) (1.8,0.44) (3.6,1.92) (5.4,2.59) (7.2,0.91)};
\foreach \p in {(0,0), (1.8,0.44), (3.6,1.92),
                (5.4,2.59), (7.2,0.91)} {
  \fill[red!75!black] \p circle (2.2pt);
}

% Best average callout (Avg. 라벨을 callout에 통합)
\node[red!85!black, font=\small, anchor=south,
      fill=white, draw=red!50, line width=0.3pt,
      rounded corners=1pt, inner sep=1.8pt]
  at (5.4,2.78) {\textbf{AVG}\ $+2.59$ pp};
% Baseline note
\node[gray!70!black, font=\small, anchor=south west]
  at (0.1,0.3) {$\lambda{=}0$ baseline};

% Axes
\draw[->, thick] (0,0) -- (7.50,0);
\draw[->, thick] (0,0) -- (0,4.60);

% X-axis title
\node[font=\small] at (3.60,-0.55)
  {anchor strength $\lambda$};
% Y-axis title
\node[font=\small, rotate=90, anchor=south] at (-0.55,2.30)
  {$\Delta$ Acc.\ vs.\ CO (pp)};
\end{tikzpicture}%
}

\vspace{0.45em}

% ============ BOTTOM: compact table ============
{\fontsize{8pt}{12pt}\selectfont
\setlength{\tabcolsep}{3pt}

\begin{tabularx}{0.97\columnwidth}{@{}l*{5}{>{\centering\arraybackslash}X}@{}}
\toprule
$\lambda$ & MMLU & SVAMP & MBPP & TydiQA & AVG \\
\midrule
$0$ (CO)
 & $45.90$
 & $36.67$
 & $34.63$
 & $60.12$
 & $44.33$ \\

$0.1$
 & \shortstack[c]{$46.16$\\[-1pt]{\scriptsize$(+0.26)$}}
 & \shortstack[c]{$37.33$\\[-1pt]{\scriptsize$(+0.66)$}}
 & \shortstack[c]{$35.02$\\[-1pt]{\scriptsize$(+0.39)$}}
 & \shortstack[c]{$60.56$\\[-1pt]{\scriptsize$(+0.44)$}}
 & \shortstack[c]{$44.77$\\[-1pt]{\scriptsize$(+0.44)$}} \\

$0.5$
 & \shortstack[c]{$46.81$\\[-1pt]{\scriptsize$(+0.91)$}}
 & \shortstack[c]{$40.00$\\[-1pt]{\scriptsize$(+3.33)$}}
 & \shortstack[c]{$37.35$\\[-1pt]{\scriptsize$(+2.72)$}}
 & \shortstack[c]{$60.85$\\[-1pt]{\scriptsize$(+0.73)$}}
 & \shortstack[c]{$46.25$\\[-1pt]{\scriptsize$(+1.92)$}} \\

\rowcolor{red!8}
$\mathbf{1.0}$
 & \shortstack[c]{$\mathbf{47.71}$\\[-1pt]{\scriptsize$(+1.81)$}}
 & \shortstack[c]{$\mathbf{40.33}$\\[-1pt]{\scriptsize$(+3.66)$}}
 & \shortstack[c]{$\mathbf{38.52}$\\[-1pt]{\scriptsize$\mathbf{(+3.89)}$}}
 & \shortstack[c]{$\mathbf{61.13}$\\[-1pt]{\scriptsize$(+1.01)$}}
 & \shortstack[c]{$\mathbf{46.92}$\\[-1pt]{\scriptsize$\mathbf{(+2.59)}$}} \\

$5.0$
 & \shortstack[c]{$46.48$\\[-1pt]{\scriptsize$(+0.58)$}}
 & \shortstack[c]{$37.67$\\[-1pt]{\scriptsize$(+1.00)$}}
 & \shortstack[c]{$36.19$\\[-1pt]{\scriptsize$(+1.56)$}}
 & \shortstack[c]{$60.62$\\[-1pt]{\scriptsize$(+0.50)$}}
 & \shortstack[c]{$45.24$\\[-1pt]{\scriptsize$(+0.91)$}} \\
\bottomrule
\end{tabularx}
}

\caption{\textbf{$\lambda$-ablation of the anchor-modulation
strength.}
LLaMA-2-7B + Alpaca-GPT4 at $\rho{=}10\%$. Gains are reported over the
CO baseline, obtained by disabling representation-anchor modulation
($\lambda{=}0$). The sweep shows an inverted-U trend, with
$\lambda{=}1.0$ serving as our empirical default.}
\label{fig:lambda-ablation}
\end{figure}

\subsection{$\lambda$-ablation of the anchor-modulation strength}
\label{sec:lambda-ablation}
\paragraph{Inverted-U effect of anchor modulation.}
Figure~\ref{fig:lambda-ablation} shows an inverted-U effect. When
\(\lambda{=}0\), the method removes representation-anchor modulation
and recovers the CO baseline to within noise. As \(\lambda\) grows,
high-projection candidates exert increasing influence on the carrier
ranking; beyond \(\lambda{=}1\), the gains decline. This pattern
supports using the projection as a bounded auxiliary signal rather
than allowing it to dominate the carrier.

\paragraph{Default choice of \(\lambda{=}1.0\).}
The sweep peaks at \(\lambda{=}1.0\), and this optimum is stable across
tasks and backbones. We therefore use \(\lambda{=}1.0\) as the
empirical default and expose only \(\rho\) in the remaining
experiments.



% ============== Table 4: Efficiency and dependency comparison ==============
% No package required for the symbols below.
\providecommand{\nodep}{--}
\providecommand{\paiddep}{\textbf{\$}}
\providecommand{\seeddep}{\textsc{Seed}}
\providecommand{\humandep}{\textsc{Human}}
\providecommand{\notmeasured}{\textsc{N/M}$^{\dagger}$}
\providecommand{\costlb}[1]{$>#1^{\dagger}$}

\begin{table}[t]
\centering
\caption{\textbf{Cost and dependency comparison} at
$\rho{=}10\%$ on LLaMA-2-7B + Alpaca-GPT4.}
\label{tab:efficiency}

{\fontsize{8pt}{11pt}\selectfont
\setlength{\tabcolsep}{3pt}

\begin{tabularx}{\columnwidth}{@{}Yccrrr@{}}
\toprule
Method
& \shortstack[c]{Ext.\\[-1pt]Scorer}
& \shortstack[c]{Cap.\\[-1pt]Ref.}
& \shortstack[c]{Sel.\\[-1pt]h}
& \shortstack[c]{SFT\\[-1pt]h}
& \shortstack[c]{Cost\\[-1pt]\$} \\
\midrule

\rowcolor{gray!15}
\multicolumn{6}{@{}l@{}}{\textbf{Static / precomputed selectors}}\\
LIMA       & \humandep & \nodep & \notmeasured & 0.28 & \costlb{0.33} \\
AlpaGasus  & \paiddep  & \nodep & \notmeasured & 0.63 & \costlb{0.73} \\
Q2Q        & \paiddep  & \nodep & \notmeasured & 0.62 & \costlb{0.71} \\
SelectIT   & \nodep    & \nodep & \notmeasured & 0.64 & \costlb{0.74} \\

\midrule

\rowcolor{gray!15}
\multicolumn{6}{@{}l@{}}{\textbf{Static / model-internal selector}}\\
NAIT       & \nodep & \seeddep & 1.18 & 0.66 & 2.11 \\

\midrule

\rowcolor{gray!15}
\multicolumn{6}{@{}l@{}}{\textbf{Dynamic / model-internal selectors}}\\
Data Agent (PPO) & \nodep & \nodep & 3.77 & 0.68 & 5.11 \\

\rowcolor{red!8}
\textbf{\textsc{TAG} (ours)}
& \nodep & \nodep
& \textbf{2.34}
& \textbf{0.64}
& \textbf{3.43} \\

\bottomrule
\end{tabularx}
}

\vspace{2pt}
\begin{minipage}{\columnwidth}
\footnotesize
\textit{Notes.} ``--'' denotes no dependency;
\humandep/\paiddep: human/paid-API scoring;
\seeddep: human-chosen in-domain seed.
GPU cost measured on A100 80GB at effective batch $128$,
normalized at \$$1.15$/A100-hr. $^{\dagger}$~SFT cost only;
selection stage not measured in our local setup. For API-based
selectors (AlpaGasus, Q2Q), API cost depends on scorer model,
candidate count, and prompt/output length (GPT-4.1 rates:
\$$2.00$/M input, \$$8.00$/M output tokens).
\end{minipage}
\end{table}

% ============== Table 5: Backbone robustness ==============

\begin{table}[t]
\centering
\caption{\textbf{Multi-backbone robustness on Alpaca-GPT4 (0.5B--7B).}}
\label{tab:backbone-robustness}
{\fontsize{8pt}{11pt}\selectfont
\setlength{\tabcolsep}{3pt}
\begin{tabularx}{\columnwidth}{@{}Yrrrrr@{}}
\toprule
Method & MMLU & SVAMP & MBPP & TydiQA & AVG \\
\midrule
\rowcolor{gray!15}
\multicolumn{6}{@{}l@{}}{\textbf{Mistral-7B}} \\
Full FT (100\%)        & 53.13              & \underline{58.24} & \textbf{51.22}     & 50.16              & 53.19 \\
Random (10\%)          & 56.80              & 57.00              & 48.64              & \underline{53.29}  & 53.93 \\
Data Agent (PPO, 10\%) & \underline{57.38}  & 56.67              & 50.28              & \textbf{54.36}     & \underline{54.67} \\
\rowcolor{red!8}
\textbf{\textsc{TAG} ($\lambda{=}1$, 10\%)} & \textbf{57.93} & \textbf{61.33} & \underline{50.97} & 53.09 & \textbf{55.83} \\
\addlinespace[3pt]

\rowcolor{gray!15}
\multicolumn{6}{@{}l@{}}{\textbf{Qwen2.5-7B}} \\
Full FT (100\%)        & \textbf{73.93}     & \underline{83.33} & 74.31              & \underline{55.78}  & \underline{71.84} \\
Random (10\%)          & 71.92              & 83.00              & \textbf{75.48}     & 54.10              & 71.13 \\
Data Agent (PPO, 10\%) & 73.59              & 83.01              & \textbf{75.48}     & 53.66              & 71.43 \\
\rowcolor{red!8}
\textbf{\textsc{TAG} ($\lambda{=}1$, 10\%)} & \underline{73.68} & \textbf{83.66} & \underline{75.10} & \textbf{56.83} & \textbf{72.32} \\
\addlinespace[3pt]

\rowcolor{gray!15}
\multicolumn{6}{@{}l@{}}{\textbf{Qwen2.5-0.5B}} \\
Full FT (100\%)        & \textbf{47.85}     & 41.33              & 41.63              & 40.27              & 42.77 \\
Random (10\%)          & 47.71              & 40.68              & 43.19              & 39.88              & 42.87 \\
Data Agent (PPO, 10\%) & 47.46              & \underline{44.00} & \underline{46.30} & \underline{41.04} & \underline{44.70} \\
\rowcolor{red!8}
\textbf{\textsc{TAG} ($\lambda{=}1$, 10\%)} & \underline{47.77} & \textbf{45.33} & \textbf{47.08} & \textbf{42.73} & \textbf{45.73} \\
\bottomrule
\end{tabularx}
}
\end{table}

% ============== Table 6: Instruction-dataset transfer (Evol-Instruct) ==============
\begin{table}[t]
\centering
\caption{\textbf{Instruction-dataset transfer on Evol-Instruct} at $\rho{=}10\%$.  Setup mirrors Table~\ref{tab:backbone-robustness}}
\label{tab:dataset-transfer}
{\fontsize{8pt}{12pt}\selectfont
\setlength{\tabcolsep}{3pt}
\begin{tabularx}{\columnwidth}{@{}Yrrrrr@{}}
\toprule
Method & MMLU & SVAMP & MBPP & TydiQA & AVG \\
\midrule
\rowcolor{gray!15}
\multicolumn{6}{@{}l@{}}{\textbf{LLaMA-2-7B}} \\
Full FT (100\%)        & \underline{45.63} & 37.67 & 35.41 & \underline{64.10} & 45.70 \\
Random (10\%)          & 45.50 & \textbf{38.67} & \textbf{38.52} & 64.00 & \underline{46.67} \\
Data Agent (PPO, 10\%) & 44.89 & 38.27 & \underline{38.44} & 63.65 & 46.31 \\
\rowcolor{red!8}
\textbf{\textsc{TAG} ($\lambda{=}1$, 10\%)} & \textbf{46.25} & \underline{38.45} & 37.69 & \textbf{64.96} & \textbf{46.84} \\
\addlinespace[3pt]
\rowcolor{gray!15}
\multicolumn{6}{@{}l@{}}{\textbf{Qwen2.5-7B}} \\
Full FT (100\%)        & \textbf{71.79} & 81.00 & 75.88 & 57.80 & 71.62 \\
Random (10\%)          & \underline{71.71} & \underline{83.33} & 75.88 & \textbf{59.10} & \textbf{72.51} \\
Data Agent (PPO, 10\%) & 71.02 & 82.74 & \textbf{76.45} & \underline{58.79} & 72.25 \\
\rowcolor{red!8}
\textbf{\textsc{TAG} ($\lambda{=}1$, 10\%)} & 71.68 & \textbf{84.16} & \underline{76.02} & 58.04 & \underline{72.48} \\
    \bottomrule
\end{tabularx}
}
\end{table}

\subsection{Cost of reference-free dynamic selection}
\label{sec:cost}
Table~\ref{tab:efficiency} reports measured local costs and required
dependencies under the same selection budget. ``Cap. ref.'' denotes
a capability reference used to guide selection, such as an in-domain
seed set manually chosen to represent the target capability space.
Static/precomputed selectors are included only as lower-bound
references, since their reported costs do not fully capture offline
dependencies such as external scoring, subset construction, or
curated references.

Among measured model-internal selectors, NAIT is cheaper than
\textsc{TAG}, but it relies on a human-chosen in-domain seed as a
capability reference. This is a strong manual prior: it requires
target-domain knowledge before selection and can bias the selected
subset toward the seed distribution. NAIT's low computational cost
is thus partly obtained by shifting difficulty to manual reference
construction, which is hard to scale, audit, and reproduce.

\textsc{TAG} removes this seed-reference dependency. It is moderately
more expensive than NAIT because state-conditioned dynamic selection
requires additional scoring, but remains cheaper than Data Agent (PPO).
The reason is structural: the carrier and anchor projections reuse
the same scoring pass, so \textsc{TAG} adds lightweight geometric
post-processing rather than a separate policy-optimization loop. Thus,
\textsc{TAG} trades moderate overhead for reference-free dynamic
selection with lower measured cost than the PPO-based selector in our
local setup.


\subsection{Backbone robustness check}
\label{sec:qwen-scale}

Table~\ref{tab:backbone-robustness} evaluates the same \textsc{TAG}
configuration (\(\lambda{=}1,\rho{=}10\%\)) on Alpaca-GPT4 across
Mistral and Qwen backbones spanning 0.5B--7B parameters; AVG averages
MMLU, SVAMP, MBPP, and TyDiQA. \textsc{TAG} achieves the best AVG in
all three backbone groups, suggesting that representation-anchor
modulation is not specific to the LLaMA-2-7B setting.


\subsection{Instruction-dataset transfer}
\label{sec:dataset-transfer}

Table~\ref{tab:dataset-transfer} replaces Alpaca-GPT4 with
WizardLM Evol-Instruct to test transfer across instruction-data
distributions. Evol-Instruct contains harder and more structured
instructions, making random \(10\%\) subsets strong and leaving
less headroom for selection gains.

\textsc{TAG} achieves the best AVG on LLaMA-2-7B. On Qwen2.5-7B,
it nearly ties the best Random baseline while still outperforming
Full FT and Data Agent. Thus, even under a stronger source-data
regime, the same representation-anchor rule remains competitive
without retuning.

